\documentclass[10pt,a4paper]{amsart}
\usepackage[margin=19mm]{geometry}
\usepackage{amsmath,amssymb,mathtools}
\usepackage[scr=rsfs,cal=euler]{mathalfa}
\usepackage{xfrac}
\usepackage[unicode,hidelinks,bookmarks=false]{hyperref}
\newcommand{\ie}{i.e.\ }
\newcommand{\cf}{cf.\ }
\newcommand{\resp}{respectively\ }
\newcommand{\Subsection}{\subsection}
\providecommand{\nobreakdash}{\nobreak\hbox{-}}
\setlength{\emergencystretch}{2em}
\multlinegap0pt
\usepackage{amssymb}
\usepackage{mathtools}
\usepackage{tikz}
\usepackage{natbib}
\usepackage[shortlabels]{enumitem}
\newtheorem{theorem}{Theorem}
\newtheorem{claim}{Claim}
\newcommand{\bK}{\mathbb K}
\newcommand{\bN}{\mathbb N}
\newcommand{\bP}{\mathbb P}
\DeclareMathOperator{\dcl}{\mathrm{dcl}}
\newcommand{\oa}{\overline{a}}
\newcommand{\ob}{\overline{b}}
\newcommand{\oz}{\overline{z}}
\DeclareMathOperator{\tp}{\mathrm{tp}}
\title{Residually Constructible Extensions}
\author{Pietro Freni, Angus Matthews}
\date{Source: arXiv:2501.10033v6 (CC BY 4.0)}
\begin{document}
\maketitle

\noindent\emph{Source and licence.} Residually Constructible Extensions. Version arXiv:2501.10033v6 (CC BY 4.0), distributed by arXiv under the Creative Commons Attribution 4.0 International licence, http://creativecommons.org/licenses/by/4.0/, as recorded in the arXiv licence metadata for this exact version and shown on the abstract page https://arxiv.org/abs/2501.10033v6. Full licence notice: \texttt{licenses/arxiv-2501.10033-CC-BY-4.0.txt}.\par\smallskip

\noindent\textit{Editorial scope.} Counted unit: the article's theorem theorem (source0.tex lines 894--896). The article's complete original proof of it is delivered with it (source0.tex lines 897--962, one \texttt{proof} environment of 66 lines). No mathematics is written, repaired or re-derived: the statement and the proof are the article's own lines, in the article's own notation, and no retained line is substituted or re-flowed.  No further source-local context is needed: every cross-reference of every retained line resolves inside the two delivered spans.  Nothing was omitted from any retained span, so the package carries no omission record.  No retained line contains a citation, so the wrapper carries no bibliography and no bibliography entry.  No retained line contains a figure environment, an image-inclusion command or drawing code, so no image is delivered and the package declares no asset dependency.  Editorial formatting only, in addition: an AMS \texttt{amsart} preamble that restates the article's own declarations and macros used by the retained lines, namely the environments \texttt{theorem}, \texttt{claim} on separate counters, together with the standard packages those lines need.  The retained environments print in this document as Theorem 1, Claim 1 in order of appearance, whereas the article prints the same items with the numbering of its own full document.  The complete native source of the article as distributed in the arXiv e-print for this exact version is bundled unchanged as \texttt{sources/171723/source0.tex}.
\begin{theorem}
Let $T$ be a countable o-minimal theory, and let $K$ be a model such that $(K,<)$ does not contain an uncountable well order. Let $K\langle x \rangle$ be an elementary extension of $\dcl_T$-dimension $1$. Then $K\langle x \rangle$ does not contain an uncountable well order.
\end{theorem}

\begin{proof}
    Assume the reverse for contradiction. Let $(z_\alpha)_{\alpha<\aleph_1}$ be a copy of $\aleph_1$ in $K\langle x \rangle$.
    Since $T$ is countable, by extracting a cofinal subsequence, we can reduce to the case where there is some $n \in \mathbb{N}$, a countable submodel $K_0$ of $K$, and a $n+1$-ary function $f$, and a sequence of $n$-tuples $(\oa^\alpha)_{ \alpha<\aleph_1}$, such that $f$ is definable over $K_0$ and $\{f(\oa^\alpha,x) : \alpha<\aleph_1\}$ is a strictly increasing sequence.
    Without loss of generality, let $n$ be minimal such that this data exists.
    Next, let $\bP$ be a countable submodel of $K$ containing $\bP$, and such that $\tp(x/\bP) \vdash \tp(x/K)$. We can do this since the type of $x$ over $K$ is determined by its cut, and $K$ does not contain a copy of $\aleph_1$, so the cut of $x$ is determined by countably many elements of $K$. Then, for $\alpha<\aleph_1$, let $K_\alpha = \bP\langle \oa^\beta : \beta<\alpha\rangle$.
    There is an easy case we must deal with.
    \begin{claim}
        If there is some $K_\alpha$ such that the set $\{ a_1^{\beta} : \beta<\aleph_1\}$ realises uncountably many distinct types over $K_\alpha$, then we reach a contradiction.
    \end{claim}
    \begin{proof}
        For ease of notation, let us expand $\bP$ to contain $K_\alpha$. Next, we will choose a cofinal subsequence of $\{\oa^{\alpha}\}$ as follows: Suppose for some $\alpha$ we have chosen a set of indices $i_\beta : \beta<\alpha$. Then $\bigcup_{\beta<\alpha}K_{i_\beta}$ is a countable extension of $\bP$, so there is some $a_1^{\gamma}$ realising a type over $\bP$ not realised in this model, by our hypothesis. Let $i_\alpha = \gamma$. Thus, by taking this cofinal subset without loss of generality we can assume for all $\alpha$, $a_1^{\alpha}$ realises a type over $\bP$ not realised in $\bigcup_{\beta<\alpha}K_\beta$.
        Now, we will inductively construct a sequence $\ob^\alpha : \alpha<\aleph_1$ such that $f(\ob^{\alpha},x)$ is strictly increasing and for each $\alpha$, there is some $\beta$ such that $f(\ob^{\alpha},x) < f(\oa^{\beta},x)$. Suppose we have constructed $\ob_\beta : \beta<\alpha$. Then, by countability there is some $\gamma<\aleph_1$ s.t. for all $\beta<\alpha$, $f(\ob^\beta,x)<f(\oa^\gamma,x)$.
        Consider $M = \bP\langle \oa^\gamma,\oa^{\gamma+2n+2}\rangle$. This is an extension of $\bP$ of $\dcl$-dimension at most $2n$. The types $\{\tp(a_1^{\gamma+j}/\bP) : 1\leq j\leq 2n+1\}$ are setwise orthogonal over $\bP$, so any extension of $\bP$ realising all of them must have $\dcl$-dimension at least $2n+1$. Thus, there is some $j \in \{1,2\cdots 2n+1\}$ such that $\tp(a_1^{\gamma+j}/\bP)$ is not realised in $M$. Consider the first-order formula
        \[
        \phi(y,r,s) = \exists \oz. \forall w \in [r,s]. f(\oa^\gamma,w)<f(y,\oz,w)<f(\oa^{\gamma+2n+2},w)
        \]
        Note that there is some interval $[r,s]$ containing $x$ such that $K \models \phi(a_1^{\gamma+j},r,s)$. Further, by the definition of $\bP$, we can assume $[r,s]$ is $\bP$-definable.
        Since the type of $a_1^{\gamma+j}$ over $\bP$ is not realised in $M$, there is some $\bP$-definable interval $[a,b]$ containing $a_1^{\gamma+j}$ such that for any $c \in [a,b]$, $K \models \phi(c,r,s)$. Choose some $b_1^{\alpha} \in [a,b] \cap \bP$.
        Then, since $K \models \phi(b_1^{\alpha},r,s)$, we can choose $b_2^{\alpha}, \cdots b_n^{\alpha}$ such that for all $w \in [r,s]$, \[f(\oa^{\gamma},w)<f(\ob^{\alpha},w)<f(\oa^{\gamma+2n+2},w)\]
        Thus, $\ob^{\alpha}$ satisfies the induction hypotheses. But then, each $b_1^{\alpha}$ is in $\bP$.
        $\bP$ is countable, so by choosing a cofinal subsequence of the $\{b^{\alpha}\}$, we can reduce to the case where $b_1^{\alpha}$ is constant as we vary $\alpha$. We can define $f_2(\oz,w) = f(b_1^\alpha,\oz,w)$ over $\bP$, and $\{f_2(b_2^\alpha, \cdots b_n^\alpha,x) : \alpha<\aleph_1\}$ is strictly increasing, contradicting the minimality of $n$.
    \end{proof}

    Thus, we can reduce to the case where over every $K_\alpha$, the set $\{a_1^{\beta}\}$ realises at most countably many types. Let 
    \[TP_{\aleph_1}(K_\alpha)=\big\{p \in S_1(K_\alpha): |\{a_1^\beta: \beta < \aleph_1,\, p(a_1^\beta)\}|>\aleph_0\big\}\]
    be the collection of types over $K_\alpha$ that have uncountably many realizations in $\{a_1^\beta: \beta < \aleph_1\}$.     Then, consider the function
    \begin{align*}
    g : \aleph_1 &\rightarrow \{0\} \cup \mathbb{N} \cup \{\aleph_0\}\\
    \alpha &\mapsto |TP_{\aleph_1}(K_\alpha)|.
    \end{align*}
    
    $g$ is increasing, as for any $\alpha<\beta$, any type $p \in TP_{\aleph_1}(K_\alpha)$ must extend to at least one type $p'$ over $K_{\beta}$, and further it must extend to one in $TP_{\aleph_1}(K_\beta)$, by Claim 1. We note also that no two distinct types over $K_\alpha$ can extend to the same one over $K_\beta$. This implies $g$ is eventually constant. Also, we know $g$ is always at least $1$. Now we split into cases.
    \begin{claim}
        If $g$ is eventually identical to the constant function $m$, for some $m \in \mathbb{N}$, we reach a contradiction.
    \end{claim}
    \begin{proof}
        Suppose toward contradiction that $g$ eventually equals $m \in \bN$. Let $g(\alpha) = m$. Let $p_\alpha \in TP_{\aleph_1}(K_\alpha)$.
       In this case, by a pigeonhole argument, for each $\alpha<\beta < \aleph_1$, $p_\alpha$ must have a unique extension $p_\beta \in TP_{\aleph_1}(K_\beta)$. Notice however that since $\bK$ does not contain uncountable well orders, there must be $\beta$ such that for all $\gamma>\beta$, $p_\beta \vdash p_\gamma$. But since no $p_\gamma$ is isolated, we would have that for every $\gamma>\beta$,  $p_\beta(y) \vdash p_\gamma(y) \vdash y \notin \bK_\gamma$ and we would have $p_\beta (\bK)=\emptyset$, yielding a contradiction.
    \end{proof}
    \begin{claim}
        If $g$ is eventually identical to the constant function $\aleph_0$, we reach a contradiction.
    \end{claim}
    \begin{proof}
        Suppose so. By enlarging $\bP$, we may assume $g$ is always equal to $\aleph_0$. We will inductively construct an increasing sequence of countable ordinals $(C_\alpha)_{\alpha<\aleph_1}$, and a sequence of tuples $(\ob^{\alpha})_{\alpha<\aleph_1}$ such that for all $\alpha$,
        \begin{enumerate}
        \item $b_1^\alpha \in \bigcup_{\beta<\alpha} K_{C_\beta}$.
        \item For all $\beta<\alpha$,
        \[f(\oa^{C_\beta},x) < f(\ob^{\alpha},x) \quad \text{and} \quad f(\ob^{\alpha},x)<f(\oa^{C_\alpha},x).\]
        \end{enumerate}
        Note that these conditions are sufficient for $f(\ob^{\alpha},x)$ to have order type $\aleph_1$.
        Suppose that for some $\alpha$, we have done all this for every $\beta<\alpha$.
        Let $M = \bigcup_{\beta<\alpha} K_{C_\beta}$. By the assumption on $g$, $TP_{\aleph_1}(M)$ has cardinality $\aleph_0$.
        Choose $i_2>i_1\geq \sup_{\beta<\alpha}(C_\beta)$ such that the set $\{a_1^{\gamma} : i_1<\gamma<i_2\}$ realises every element of $TP_{\aleph_1}(M)$. Let the realising indices be $\delta_1,\delta_2, \cdots$. Consider the following first-order formula
        \[
        \phi(y,r,s) = \exists \oz. \forall w \in [r,s]. f(\oa^{i_1},w)<f(y,\oz,w)<f(\oa^{i_2},w)
        \]
        Note that since $x$ is not definable in $K$, for any $i_1<\gamma<i_2$, there are $r,s \in K$ such that $K \models \phi(a_1^{\gamma},r,s)$.
        Note also that by definition of $\bP$, $\tp(x/\bP) \vdash \tp(x/K)$, so we can assume $r,s \in \bP$.
        By o-minimality, there is some constant $N$ such that for any fixed $r,s$, $\phi(y,r,s)$ defines a subset of $K$ consisting of at most $N$ points and intervals.
        Choose some $r,s \in \bP$ close enough to $x$ that $K \models \phi(a_1^{\delta_l},r,s)$ for all $1\leq l \leq N+1$. Thus, since these are $N+1$ distinct points, two of them must be in the same $K$-definable interval $I$ such that $K \models \phi(I,r,s)$. But since these two have distinct cuts over $M$, $I$ must contain some point in $M$.
        We will suggestively call this point $b_1^{\alpha}$. Then, choose points $b_2^{\alpha}, \cdots b_n^{\alpha}$ in $K$, witnessing $\phi(b_1^{\alpha},r,s)$. Note that since $f(b^{\alpha},x)>f(a^{i_1},x)$, for all $\beta<\alpha$, $f(b^{\alpha},x)>f(a^{C_\beta},x)$, as desired. Then, if we set $C_\alpha=i_2$, $b^{\alpha}$ and $C_\alpha$ fulfill the inductive criteria.
        Observe that each $b_1^\alpha$ is definable over $\bigcup_{\beta<\alpha} K_{C_\beta}$. Thus, the induction works. Since definable closure is finitary, for each $\alpha>0$, there is some $h(\alpha)<\alpha$ such that $b_1^\alpha$ is definable over $K_{C_{h(\alpha)}}$. Then, since $h$ a function strictly below the identity from $\aleph_1$ to $\aleph_1$, by Fodor's lemma there is some $C_\alpha$ with uncountable preimage.
        This implies that by taking a cofinal subset of $\{\ob^\alpha : \alpha<\aleph_1\}$, we may reduce to the case where all the $b_1^{\alpha}$ are definable over $K_{C_\alpha}$. Then, by taking another cofinal subset, we may reduce to the case where all the $b_1^{\alpha}$ are equal. This contradicts the minimality of $n$.
    \end{proof}
    All cases lead to a contradiction, so the theorem is proved.
\end{proof}

\end{document}
